% Chapter 3. Material: docs/requirements/quality-scenarios.md, docs/research/metrics-gqm.md,
% docs/research/preregistration.md, docs/data/annotation-guide.md, docs/data/reports.md,
% docs/architecture/design-iterations.md, docs/risks.md. Figure R-06 → fig-harness (Ch. 4).
\chapter{Methodology}\label{ch:method}

This chapter explains how the thesis produces and evaluates its artefacts. Section~\ref{sec:method-dsr}
introduces the research paradigm. Sections~\ref{sec:method-req}--\ref{sec:method-gqm} describe how
requirements, architecture and metrics were derived. Sections~\ref{sec:method-hyp}--\ref{sec:method-s4}
present the four studies, and Sections~\ref{sec:method-prov}--\ref{sec:method-validity} cover
preregistration, ethics and threats to validity.

\section{Research paradigm}\label{sec:method-dsr}

The thesis follows design science research, which produces knowledge by building an artefact that
addresses a real problem and evaluating it rigorously \cite{SRC-100}. Wieringa separates two kinds of
work inside such a project: \emph{design problems}, which ask how to change the world (``build a tutor
that \ldots''), and \emph{knowledge questions}, which ask what is true about it (``how often does the model
reject correct reading?'') \cite{SRC-101}. The research questions of this thesis are knowledge questions;
the objectives O1--O3 and O6 are design problems that make them answerable. Table~\ref{tab:hevner} shows
how the thesis addresses Hevner's seven guidelines.

\begin{table}[ht]
\centering\small
\caption{Hevner's design science guidelines \cite{SRC-100} applied to this thesis.}
\label{tab:hevner}
\begin{tabularx}{\linewidth}{@{}P{3.6cm}L@{}}
\toprule
Guideline & In this thesis \\
\midrule
Design as an artefact & Artefact A, the reading tutor; artefact B, the development harness. \\
Problem relevance & Beginners practise alone; teachers' time does not scale; no tool checks primer-level
reading in the learners' language. \\
Design evaluation & Measurement against two teachers (RQ-1), error analysis (RQ-2), pilot experiment
(RQ-3), case study (RQ-4). \\
Research contributions & Evaluation on an unstudied population and material; dataset; method. \\
Research rigour & Preregistration; standard MDD metrics; agreement statistics with confidence intervals;
case-study protocol \cite{SRC-102}; Empirical Standards as checklist \cite{SRC-103}. \\
Design as a search process & 50 recorded decisions with rejected alternatives; seven design iterations. \\
Communication & Thesis, dataset card, open architecture, and a doctoral roadmap. \\
\bottomrule
\end{tabularx}
\end{table}

\section{Requirements}\label{sec:method-req}

\subsection{The audience as the main constraint}

One constraint was stated first and checked against every later decision: the learners are
Russian-speaking adults who learn to read Arabic from zero. Seven rules follow from it: explain every sign
in Russian before it is needed; follow the book's order in small steps; let the learner hear before
reading; explain errors through familiar sounds and show them in the text; consolidate with short
questions; never punish; and treat the Qur'anic text with respect, using the exact text of the Medina
mushaf and starting with \emph{isti'adha} and \emph{basmala}.

\subsection{Functional requirements and user stories}

Requirements were collected from the author's requests, the prototype, the book and the literature, and
written with stable identifiers and MoSCoW priorities. The current specification has 47 functional
requirements, grouped as follows: curriculum from the book (FR-01, FR-16, FR-36--FR-41), assessment and
feedback (FR-03--FR-05, FR-20--FR-22, FR-43), learner experience (FR-02, FR-07, FR-31--FR-35), research
data (FR-06, FR-11--FR-14, FR-17, FR-24, FR-26--FR-29), privacy (FR-08--FR-10, FR-44) and
later interface additions (FR-45--FR-47). User stories
describe the same requirements from the viewpoints of learner, teacher and researcher. A traceability matrix
links each requirement to code and tests; an automated report shows which requirements are covered by
passing tests (59 of 60 at the time of writing).

\subsection{Quality attribute scenarios}

Non-functional requirements were rewritten as six-part quality attribute scenarios --- source, stimulus,
environment, artefact, response and response measure --- following Bass, Clements and Kazman
\cite{SRC-047}. Scenarios were prioritised in a utility tree by importance and technical risk.
Table~\ref{tab:qas} lists them.

\begin{table}[ht]
\centering\small
\caption{Quality attribute scenarios (QAS-01--QAS-11).}
\label{tab:qas}
\begin{tabularx}{\linewidth}{@{}lP{2.6cm}L@{}}
\toprule
ID & Attribute & Scenario and measure \\
\midrule
QAS-01 & Accuracy & The verdict agrees with a teacher: precision, recall, F1; $\kappa$ bot--teacher vs.\
teacher--teacher. \\
QAS-02 & Accuracy & Correct reference reading is accepted: share of failed verdicts on reference audio
$\le$5\,\% (to approve). \\
QAS-03 & Performance & Response to a voice message: 90\,\% of attempts processed in $\le$10\,s. \\
QAS-04 & Reliability & A model failure does not break the lesson; 100\,\% of failures logged and
answered. \\
QAS-05 & Reliability & Position in the book survives a restart: no lost positions. \\
QAS-06 & Input robustness & Too short, silent or over-long recordings are not graded. \\
QAS-07 & Usability & At least four of five beginners reach a passed line in their first session without
help (to approve). \\
QAS-08 & Privacy & No audio without consent; deletion leaves no rows or files. \\
QAS-09 & Reproducibility & Every attempt can be recomputed: versions, thresholds and hash stored. \\
QAS-10 & Modifiability & A new recogniser is added without touching the domain. \\
QAS-11 & Process integrity & The agent cannot weaken tests during implementation. \\
\bottomrule
\end{tabularx}
\end{table}

\section{Architecture design method}\label{sec:method-add}

The architecture was derived with a lightweight form of Attribute-Driven Design \cite{SRC-047}: each
iteration takes one goal and a few drivers (scenarios, requirements, constraints), records decisions as
architecture decision records (ADRs) \cite{SRC-010}, and closes with a test that confirms the decision.
For a single author full iteration documents would be bureaucracy, so an iteration is one line in a table
and the details live in the ADRs, including rejected options. Seven iterations were completed: replaceable
model, database and messenger; research data with personal data kept apart; trustworthy and contestable
assessment; curriculum strictly from the book; an interface usable without a teacher; agent development
that cannot be gamed; and diagnostics without leaking personal data. Chapter~\ref{ch:impl} describes the
results.

\section{Measurement framework}\label{sec:method-gqm}

Metrics were derived with the Goal--Question--Metric method \cite{SRC-048}: a goal is stated in a fixed
template (analyse \emph{object} for the purpose of \emph{purpose} with respect to \emph{quality focus} from
the viewpoint of \emph{stakeholder} in the context of \emph{context}), refined into questions, and only then
into metrics and data. A metric without a goal is not collected. Six goals were defined: G1 agreement with
teachers (RQ-1), G2 beginners' errors (RQ-2), G3 effect of explanations (RQ-3), G4 engagement and the
beginner's path, G5 system quality, and G6 development with AI agents (RQ-4).

For the pilot, one \emph{main metric} and one \emph{counter-metric} were chosen. The main metric is the
number of book lines passed by learners for the first time in a week: it grows only when learners move
through the book. Because it can be inflated by lowering the pass threshold, the counter-metric is the
share of falsely accepted attempts, attempts passed by the tutor but not by the teacher's transcription.

\section{Hypotheses}\label{sec:method-hyp}

Table~\ref{tab:hyp} lists the confirmatory hypotheses. They are recorded in the preregistration
(Section~\ref{sec:method-prov}); every other analysis is labelled exploratory and its findings are offered
as hypotheses for future work.

\begin{table}[ht]
\centering\small
\caption{Confirmatory hypotheses (preregistration v0.2).}
\label{tab:hyp}
\begin{tabularx}{\linewidth}{@{}llL@{}}
\toprule
ID & RQ & Hypothesis \\
\midrule
H1.1 & RQ-1 & Precision and recall of error detection on beginners are lower than reported on the model's
training domain. \\
H1.2 & RQ-1 & False rejections (FRR) are more frequent than false acceptances (FAR). \\
H1.3 & RQ-1 & $\kappa$ between tutor and teacher is lower than $\kappa$ between teachers. \\
H2.1 & RQ-2 & Errors on \ar{ق} and \ar{غ} are less frequent among Kazakh speakers than among Russian
speakers (tested only if both groups are present). \\
H2.2 & RQ-2 & Russian speakers make more errors in emphasis than in place of articulation. \\
H3.1 & RQ-3 & The score gain on the next attempt is larger after a full explanation than after a
correct/incorrect signal. \\
H3.2 & RQ-3 & The difference in H3.1 persists on a delayed test after 1--2 weeks. \\
\bottomrule
\end{tabularx}
\end{table}

\section{Participants and materials}\label{sec:method-participants}

The pilot is closed: learners are invited through teachers and the author's acquaintances, and the target is
20--40 adult learners over 4--6 weeks. Participation in research requires a separate consent; a short entry
questionnaire, asked only of consenting learners, records age (18+), first language and prior experience of
reading Arabic. Learners who do not consent can still use the tutor; their attempts are graded but no audio
is kept and they are excluded from the analysis. The materials are the cards of the letter and vowel lessons
(pages 10--24 of the book) and, for exploratory analysis, later lessons that learners reach. The unit of
analysis is an attempt; learners are clusters, because attempts by one learner are not independent.

Exclusions are fixed in advance: attempts graded by the development stub; recordings labelled unusable or
off-task by an annotator (counted separately); and attempts before the preregistration freeze.

\section{Study 1: agreement with teachers (RQ-1)}\label{sec:method-s1}

\paragraph{Annotation} Two teachers who read confidently with tajwid annotate every attempt. Each sees the
line of the book, the recording, the canonical phonemes, the phonemes the tutor heard and the errors it
reported. If the tutor heard correctly, the annotator marks \emph{agree}; otherwise \emph{fix} and enters the
phonemes that were actually pronounced. Two further marks are possible: \emph{audio unusable} and
\emph{off task}. A separate flag marks a \emph{self-correction}, when the learner made an error and immediately
re-read correctly in the same message; following the taxonomy of Al Mdfaa et al.\ \cite{SRC-068}, these
attempts are reported in their own block and excluded from the main RQ-1 figures. Annotation can be done in
Telegram or in a spreadsheet package; both write to the same table with the annotator code and the version
of the annotation guide. Annotators see the tutor's answer; this speeds up work but may bias them towards
agreement, which is treated as a threat to validity.

\paragraph{Phoneme-level metrics} For each attempt the canonical, annotated and recognised phoneme sequences
are aligned, and each canonical phoneme is classified as TA, TR, FA or FR (Section~\ref{sec:lr-mdd}).
Precision, recall, F1 and diagnosis accuracy are computed over all phonemes, by error type and by letter
\cite{SRC-064}.

\paragraph{Verdict-level metrics} The teacher's verdict is obtained by applying the same grading policy to
the teacher's transcription, with the thresholds stored in the attempt. FAR is the share of attempts the
tutor passed although the teacher's transcription fails; FRR is the share it failed although the
transcription passes. Agreement is measured with Cohen's $\kappa$ \cite{SRC-104},
\begin{equation}
  \kappa = \frac{p_o - p_e}{1 - p_e},
\end{equation}
where $p_o$ is observed and $p_e$ chance agreement, computed between the tutor and each teacher and between
the two teachers. The teacher--teacher value is the ceiling against which the tutor is judged (H1.3).
Agreement between annotators on individual phonemes is measured with Krippendorff's $\alpha$ for nominal
data, which handles missing labels and any number of coders \cite{SRC-065}.

\paragraph{Uncertainty} 95\,\% confidence intervals are obtained by a cluster bootstrap that resamples
learners with replacement. Comparisons with published benchmarks are qualitative only, because the data
differ.

\section{Study 2: beginners' errors (RQ-2)}\label{sec:method-s2}

Error frequencies by type, confusion matrices ``expected $\to$ heard'' and difficulty by book line and by
the book's minimal pairs are computed twice: over all errors reported by the tutor and over errors
confirmed by annotators (true rejections). The difference between the two shows how much of an apparent
learner profile is produced by the model. To avoid conclusions driven by one active learner, every
frequency is reported with the number of learners who contribute to it. H2.1 is tested with a mixed
logistic model, with first language as a fixed effect and learner as a random effect, if both groups are
large enough; otherwise it is reported descriptively. H2.2 uses the attribute (\emph{sifat}) outputs of the
model and is reported descriptively with intervals.

\section{Study 3: effect of explanations (RQ-3)}\label{sec:method-s3}

\paragraph{Design} With 20--40 learners, two independent groups would be underpowered: a medium effect
($d\approx0.6$) would need about 45 learners per group at $\alpha=0.05$ and 80\,\% power. The study
therefore uses a within-subject design, following Neri et al.\ \cite{SRC-054}: for each learner, a random
half of the cards gives a full explanation after a failed attempt, and the other half only a
correct/incorrect signal. Assignment is fixed per learner and card, so a learner always receives the same
condition on the same card.

\paragraph{Outcomes} The primary outcome is the change in score from a failed attempt to the next attempt on
the same card. Secondary outcomes are whether the error named in the explanation recurs at the same
position, and the score on the same cards in a delayed test after 1--2 weeks.

\paragraph{Analysis} The effect is estimated with a linear mixed model
\begin{equation}
  \Delta s_{ij} = \beta_0 + \beta_1\,\mathrm{Full}_{ij} + u_i + v_j + \varepsilon_{ij},
\end{equation}
where $\Delta s_{ij}$ is the score change of learner $i$ on card $j$, $\mathrm{Full}_{ij}$ indicates the
full-explanation condition, and $u_i$, $v_j$ are random intercepts for learner and card. H3.2 uses the same
model on delayed-test scores. Across the seven confirmatory hypotheses, the Holm procedure controls the
family-wise error rate at $\alpha=0.05$; effect sizes and confidence intervals are always reported. Data
collection stops at the planned end of the pilot, not when results become significant; interim reports
are descriptive only.

\section{Study 4: the development harness (RQ-4)}\label{sec:method-s4}

The harness is evaluated as a single embedded case study \cite{SRC-102}. The case is the development of the
tutor; the units of analysis are working sessions and individual changes. Four data sources are used: the
hook journal, which records every block and approval request with time, hook and reason; session summaries,
which record requests, decisions, deviations and unverified items; the log of every prompt to the agent; and
the version history with its ADRs. Measures follow goal G6: blocks and approval requests by reason;
human interventions per session (corrections, stop signals, rejected changes); iterations until the checks
pass; and defects that passed all checks and were found later. Incidents described in the journals are
coded thematically by failure mode (test gaming, scope creep, unverified claims, data exposure,
fabricated dependencies). Self-reported productivity is deliberately not used \cite{SRC-030}.

\section{Preregistration, provenance and reproducibility}\label{sec:method-prov}

Three mechanisms protect the results from being shaped after the fact.

\paragraph{Preregistration} Hypotheses, exclusions, sample size and analysis are written down before data
collection. When the author and supervisor approve the plan, its SHA-256 fingerprint and date are recorded,
and any later change requires a new version with a stated reason. Freezing data and evaluator by hash
follows the practice of the supervisor's group \cite{SRC-067}.

\paragraph{Protected evaluator} The module that computes RQ-1 metrics is a human-owned zone: the coding agent
must ask before changing it, because an agent that can edit the evaluator can improve the score without
improving the system \cite{SRC-027,SRC-082}.

\paragraph{Provenance} Every attempt stores the canonical variants, the recognised phonemes, the raw model
output, the thresholds, the code, curriculum and model versions, the book line as shown to the learner, and
the SHA-256 of the audio. A traceability command prints fingerprints of code, curriculum, SQL reports,
migrations, dependencies and preregistration. A number enters the thesis only if the corresponding analysis
is implemented and the figure or table has a passport: report, date, number of attempts and learners, and
code fingerprint.

\section{Ethics}\label{sec:method-ethics}

The study involves voice recordings on a religious topic, so it is designed as if the strictest reading of
the law applied. Consent to the service and consent to voice recording are separate; audio is never stored
without consent; consent can be withdrawn, which deletes recordings and their annotations; every learner can
delete all data with one button; only pseudonyms appear in reports and datasets; research participation is
limited to adults; no voice prints or speaker embeddings are computed; annotators see only the line and the
tutor's answer, never the learner's identity, and recordings are sent to them in protected mode
\cite{SRC-033,SRC-036,SRC-037}. The questions of who is the data operator, whether an ethics committee review
is required and where the server is located are institutional and are being settled with the supervisor
before learners outside the closed circle are invited.

\section{Threats to validity}\label{sec:method-validity}

Table~\ref{tab:validity} lists the main threats in the classification of Wohlin et al.\ \cite{SRC-106} and
how the design addresses them.

\begin{table}[ht]
\centering\small
\caption{Threats to validity and responses.}
\label{tab:validity}
\begin{tabularx}{\linewidth}{@{}P{2.2cm}LL@{}}
\toprule
Type & Threat & Response \\
\midrule
Construct & An annotator's label is not the ``true'' error; teachers may disagree. & Two annotators;
teacher--teacher $\kappa$ as ceiling; written guide with versions. \\
Construct & Annotators see the tutor's answer and may agree by suggestion. & Reported as a limitation;
agreement analysed separately for passed and failed attempts. \\
Internal & Improvement on a second attempt may be repetition, not explanation. & Within-subject control
condition (RQ-3). \\
Internal & Long lines dilute errors in the similarity score. & RQ-1 is computed per phoneme, not only per
verdict. \\
Internal & The agent may game its own checks. & Frozen tests, protected evaluator, logged blocks
(RQ-4). \\
External & Closed pilot of acquaintances; small sample. & Described as a pilot; effect sizes with
intervals; no generalisation beyond the population. \\
External & The model was trained on the Qur'an; pages 10--37 are not Qur'anic. & Results reported by
material type. \\
Conclusion & Dependent attempts; multiple tests. & Cluster bootstrap; mixed models; Holm correction. \\
\bottomrule
\end{tabularx}
\end{table}
